% from ega4-9-fr.tex
\oldpage[IV-3]{78}

Soit $(p,q)$ un couple d'entiers $>0$ tels que $p+q=d$~; introduisons des
indéterminées $\mathrm T'_\beta$, $\mathrm T''_\gamma$ pour tous les systèmes
d'entiers $\beta$, $\gamma$ tels que $|\beta|\leqslant p$ et
$|\gamma|\leqslant q$~; pour tout système d'entiers $\alpha$ tel que
$|\alpha|\leqslant d$, considérons le polynôme de
\[
B=A[\mathrm T'_\beta,\mathrm T''_\gamma]_{|\beta|\leqslant p,
  |\gamma|\leqslant q}:
\qquad
P_\alpha(\mathrm T'_\beta,\mathrm T''_\gamma)
=\sum_{\beta+\gamma=\alpha}\mathrm T'_\beta\mathrm T''_\gamma-c_\alpha.
\]
Soit $\Omega$ une clôture algébrique de $K$~; dire qu'il existe deux polynômes
$F_1$, $F_2$ de $\Omega[\mathrm T_1,\ldots,\mathrm T_n]$, de degrés respectifs
$p$ et $q$, tels que $F_1F_2=F_{(\Omega)}$, signifie que le système d'équations
$P_\alpha(t'_\beta,t''_\gamma)=0$ ($|\alpha|\leqslant d$) admet une solution
$(t'_\beta,t''_\gamma)$ ($|\beta|\leqslant p$, $|\gamma|\leqslant q$) formée
d'éléments de $\Omega$. Soit $\mathfrak a$ l'idéal de $B$ engendré par les
$P_\alpha$~; l'interprétation précédente, et le théorème des zéros de Hilbert,
montrent que l'hypothèse sur $F_{(K)}$ implique que
$V(\mathfrak a_\eta)=\varnothing$, en désignant par $\eta$ le point générique de
$\operatorname{Spec}(A)$~; le lemme (9.7.5.1) prouve donc que dans un voisinage
de $\eta$, on a $V(\mathfrak a_x)=0$, et par suite pour ces valeurs de $x$,
$F_{(\boldsymbol{k}(x))}$ n'admet pas de factorisation
$F_{(\boldsymbol{k}(x))}=G_1G_2$, où $G_1$, $G_2$ sont des polynômes de degrés
respectifs $p$ et $q$, dont les coefficients sont dans une clôture algébrique
de $\boldsymbol{k}(x)$. Il suffit d'appliquer ce résultat à tous les couples
d'entiers $(p,q)$ tels que $p>0$, $q>0$ et $p+q=d$, pour obtenir la conclusion
du lemme (9.7.5).

\begin{proposition}[9.7.6]
\label{IV.9.7.6-fr}
Soient $S$ un préschéma noethérien intègre de point générique $\eta$, $f:X\to S$
un morphisme de type fini, $\mathcal F$ un $\mathcal O_X$-Module cohérent. Si
$\mathcal F_\eta$ est sans cycle premier associé immergé, il existe un voisinage
$U$ de $\eta$ dans $S$ tel que, pour tout $s\in U$, $\mathcal F_s$ soit sans cycle
premier associé immergé.
\end{proposition}

Comme $X_\eta$ est noethérien, il existe un homomorphisme injectif
$v:\mathcal F_\eta\to\bigoplus_{i=1}^m\mathcal G_i$, où chacun des $\mathcal G_i$ est
irrédondant et $\operatorname{Ass}(\mathcal F_\eta)$ est l'ensemble des $x_i$, où
$\{x_i\}=\operatorname{Ass}(\mathcal G_i)$ (3.2.6)~; si $Z_i$ est l'adhérence de
$\{x_i\}$ dans $X_\eta$, on a $Z_i=\operatorname{Supp}(\mathcal G_i)$ (3.1.4), et
par hypothèse les $Z_i$ sont les composantes irréductibles de
$Z=\operatorname{Supp}(\mathcal F_\eta)$. Il résulte de (8.5.2, (i) et (ii)) (en
restreignant au besoin $S$ à un voisinage de $\eta$) et (8.5.8) qu'il existe des
$\mathcal O_X$-Modules cohérents $\mathcal F_i$ et un homomorphisme injectif
$u:\mathcal F\to\bigoplus_{i=1}^m\mathcal F_i$ tels que $(\mathcal F_i)_\eta=\mathcal G_i$ pour
tout $i$ et $v=u_\eta$. Si $Y_i=\operatorname{Supp}(\mathcal F_i)$, on a
$(Y_i)_s=\operatorname{Supp}((\mathcal F_i)_s)$ pour tout $s\in S$
(\textbf{I}, 9.1.13) et en particulier $(Y_i)_\eta=Z_i$ pour tout $i$.
L'hypothèse entraîne que pour $i\neq j$, $Z_i-(Z_i\cap Z_j)$ est ouvert dense
dans $Z_i$~; on déduit donc de (9.5.3) et (9.5.4) que dans un voisinage de
$\eta$, $(Y_i)_s-((Y_i)_s\cap(Y_j)_s)$ est ouvert et dense dans $(Y_i)_s$.
Supposons que l'on ait prouvé que chacun des $(\mathcal F_i)_s$ est sans cycle
premier associé pour $s\in U$. Alors les éléments de
$\operatorname{Ass}((\mathcal F_i)_s)$ sont les points maximaux de $(Y_i)_s$~;
aucun d'eux ne peut donc appartenir à un $(Y_j)_s$ pour $i\neq j$, et la
proposition sera démontrée. On peut donc supposer que $\mathcal F_\eta$ est
\emph{irrédondant}~; en outre, on peut se borner au cas où
$S=\operatorname{Spec}(A)$ est affine~; $X$ est alors réunion d'un nombre fini
d'ouverts affines $V_j$, et si $V_j\cap\operatorname{Supp}(\mathcal F_\eta)
=\varnothing$, on aura aussi $V_j\cap\operatorname{Supp}(\mathcal F_s)=\varnothing$
pour $s$ voisin de $\eta$ (9.5.1), donc on peut se borner au cas où
$X=\operatorname{Spec}(B)$ est aussi affine et où le morphisme $f:X\to S$ est
dominant~; $A$ est donc un anneau noethérien intègre, \emph{sous-anneau} d'une
$A$-algèbre de type fini $B$, et $\mathcal F=\widetilde M$, où $M$ est un
$B$-module de type fini~; par hypothèse, si $K$ est le corps des fractions de
$A$, le $B_{(K)}$-module $M_{(K)}$ est irrédondant. Soit $\mathfrak q$ l'unique

